\section{Four oriented power blocks and three integer profiles}
\label{sec:pooling}

The signed growth interface of Section~\ref{sec:signed} composes with the
oriented-pooling theorem of~\cite{signedhorizons}. The standalone growth
certificate leaves fifteen ordinary power requests at four bases:
\(B,D,q,q+1\), with respective multiplicities \(1,1,3,10\).
Pooling their private certificates and performing checked unit-pivot
substitutions gives the following complete polynomial circuits over
unrestricted integers.

\begin{center}
\begin{tabular}{rrrrr}
\toprule
Variables & Exact degree & Residual squares & Binary operations & Height bound\\
\midrule
130 & 12 & 70 & 2475 & \(<2^{54}\)\\
116 & 16 & 56 & 2433 & \(<2^{84}\)\\
107 & 32 & 47 & 2419 & \(<2^{107}\)\\
\bottomrule
\end{tabular}
\end{center}
Binary operations count additions, subtractions, and multiplications in the
specified circuit. The last column bounds the absolute coefficients after
full expansion; these are rigorous upper bounds, not measured exact heights.
The largest literal integer constant is \(2^{20}\). The first two profiles
use the staged partition described below, with different elimination lists;
the third uses one four-base pool. Each saves one variable at its matched
degree relative to the inspected \(131,117,108\)-variable predecessors.
The degree-16 and degree-32 predecessors use smaller literal constants,
so the new forms do not improve every size measure.

\subsection{The exact identity and the strict rational-error argument}

Use \(A\) for the common Pell center, distinct from the standalone growth
center \(4t+1\). The retained frame and anchor supply
\[
 W\ge258,\qquad \abs A>W^W,\qquad
 n^2,m_i^2,k_i^2<W,
\]
and the canonical signed Pell pair \((X_i,Y_i)\) at each requested index
\(k_i\). Here \(n\) is the actual base and \(m_i\) its proposed power
output. Set \(U_i=X_i-AY_i\), using the retained multiplication wire,
and, within a block of common base \(n\), define
\begin{align}
 T_n&=2An-n^2-1,&
 H_n&=\sum_iY_i(m_i-X_i),\label{eq:pool-T-H}\\
 E_n&=\sum_i\bigl((m_i-U_i)^2-Y_i^2\bigr)-2nH_n,&
 \delta_i&=m_i-U_i-nY_i.\label{eq:pool-E}
\end{align}
Expansion gives the exact polynomial identity
\begin{equation}\label{eq:pool-identity}
 E_n-T_n\sum_iY_i^2=\sum_i\delta_i^2.
\end{equation}
For any selected blocks \(J\), the weighted pool is the single residual
\begin{equation}\label{eq:weighted-pool}
 R_J=\sum_{j\in J}n_jE_{n_j}
              \prod_{\ell\in J\setminus\{j\}}T_{n_\ell}
       -z_J\prod_{j\in J}T_{n_j}.
\end{equation}
All products are finite. This residual and its circuit contain no divisions
or variable exponentiation.

Here is a proof of its unbounded arithmetic meaning. Put
\(b=\lfloor\sqrt{W-1}\rfloor\ge16\). If \(n\ne0\), then
\[
 \abs{T_n}\ge2\abs A-W>W^W,\qquad
 \operatorname{sign}(n/T_n)=\operatorname{sign}(A).
\]
For \(n=0\), \(T_n=-1\), and the block's weighted contribution is
zero. Dividing \(R_J=0\) by its nonzero denominator product and using
\eqref{eq:pool-identity} makes
\(Q=\sum_{j,i}n_j\delta_i^2/T_{n_j}\) an integer.
Let \(J_*\) contain the nonzero bases and put
\[
 G=\prod_{j\in J_*}n_j^b,\qquad
 e_i=Gm_i-(G/n_j^b)n_j^{b+k_i}\quad(j\in J_*).
\]
These are integers; \(G=1\) for an empty product. Since
\(T_n\equiv-1\pmod n\), the canonical Pell congruence also permits
negative indices, interpreted by modular inverses, and yields
\(G\delta_i\equiv e_i\pmod{T_{n_j}}\). Therefore
\[
 Q'=\sum_{j\in J_*}\sum_i n_je_i^2/T_{n_j}
\]
differs from \(G^2Q\) by an integer. No coprimality between different
\(T_n\)'s is used. There are at most four blocks and fifteen requests,
so \(\abs{e_i}\le2b^{5b}\) and
\[
 \sum_{j,i}\abs{n_j}e_i^2
 \le60b^{10b+1}<W^{5b+2}\le W^W.
\]
The last step uses \(5b+2\le b^2+1\le W\). Thus \(\abs{Q'}<1\).
Its integrality forces \(Q'=0\), and the common sign of every
\(n_j/T_{n_j}\) forces every \(e_i=0\). Hence every nonzero-base
request is exact in \(\mathbb Q\): \(m_i=n_j^{k_i}\). If
\(\abs{n_j}\ge2\), integrality further forces \(k_i\ge0\).
At the unit bases \(\pm1\), negative indices remain allowed.

Conversely, exact nonzero-base powers give \(T_n\mid\delta_i\),
then \(T_n\mid E_n\), by \eqref{eq:pool-identity}. The integer
\(z_J=\sum_{j\in J_*}n_jE_{n_j}/T_{n_j}\) completes the pool.
Zero-base blocks remain deliberately unconstrained until the bootstrap
excludes them.

\subsection{The two partitions and the order of sign recovery}

The 107-variable profile uses \eqref{eq:weighted-pool} for all four
bases. The staged profiles instead use it for \(B,D\), followed by
the unweighted residual
\begin{equation}\label{eq:unweighted-pool}
 E_qT_{q+1}+E_{q+1}T_q-z_*T_qT_{q+1}.
\end{equation}
The unweighted argument is the same integer-error proof with weights one,
provided the two denominators have the same sign. Its numerator bound is
no larger than the bound above.

This condition is recovered in the required order. The literal horizon
and growth-divisibility row make \(F\) even and nonzero; \(C\) is even.
Thus \(B\eta=C+F+1\) makes \(B\ne0\), while
\(D=2(C+F+1)\equiv2\pmod4\) gives \(\abs D\ge2\).
The weighted block therefore recovers \(q=D^x\) and \(x\ge0\).
Consequently \(q\notin\{0,-1\}\): for \(x=0\) it is one, and
otherwise its absolute value is at least two. Hence \(q,q+1\) are
nonzero and have the same sign, as do \(T_q,T_{q+1}\).
Equation~\eqref{eq:unweighted-pool} is now legitimate; in the all-base
profile the two previously conditional blocks are now active. Every
retained power is exact. The endpoint argument in
Section~\ref{sec:signed} then gives \(x\ge1\), uses the literal request
\(q^t\) to force \(t\ge2\), and recovers all remaining signs.
No power at a zero base, nor a nonnegative-index conclusion at a unit
base, is used along this route.

\subsection{Integral restoration and what the finite checks establish}

Pooling removes eight private coordinates and eight rows from the pinned
216-variable source, then adds two pool quotients for the staged form or
one for the all-base form. The deleted singleton certificates restore
their modulus as \(T_n\) and their old quotient as
\(-\delta_i/T_n\). A deleted common-base pair restores as
\[
 L_n=E_n/T_n,\qquad K_n=2H_n+(2A-n)L_n.
\]
The exact-power congruence proves that these quotients are integral on
zeros. These restoration maps are not asserted to be polynomials on all
integer tuples. Conversely, every original certificate supplies the new
integer pool quotients, so the old and pooled systems have the same
projection onto retained coordinates.

The final profiles apply respectively 80, 94, and 102 acyclic unit pivots.
Each pivot has the form \(\varepsilon v+f=0\),
\(\varepsilon\in\{1,-1\}\), and is restored by the integer polynomial
\(v=-\varepsilon f\) in the certified reverse order. This step gives
actual polynomial zero-set bijections. The output circuit is the sum of
the squares of all remaining residuals.

The independent checker binds the actual 34 frame atoms, 13 principal
Pell pairs, 15 requests, historical affine offsets, private-coordinate
deletions, and every local polynomial identity and pivot. The complete
request table and restoration data are in \artifact{pooling/PROOF.md}
and \artifact{pooling/certificates/}. Nonzero homogeneous specializations
prove the exact degrees; coefficient-norm propagation proves the stated
height bounds. Both ordinary and optimized-Python verification pass,
including fifteen deliberately corrupted profiles that are rejected.
These finite checks authenticate the literal equations. The unbounded
pooling proof, the signed-horizon soundness and cofinal completeness,
and the named inherited logical calibration supply the equivalence with
\(\Con(\ZFC)\); no new end-to-end proof-assistant build is claimed.
